\section{Conclusion}
\label{sec:conclusion}

\subsection{Summary}
\label{subsec:summary}
This work presents an architecture drawing on approaches by~\cite{Kim2021} and~\cite{Kolter2020}
to optimize the parameters of multiple SABR models using a neural network with implicit deep equilibrium layers as the candidate set selection mechanism.
Our approach was successful in enforcing the no-arbitrage conditions in the resulting surfaces by solving for a fixed-point problem between the observed data and generated IV values.
The time complexity of the trained model vastly outperformed the Monte Carlo method, with a speedup of $6\,300\%$, and was able to generalize well to test data,
as shown in \Cref{tab:training} and \Cref{tab:performance}, respectively.
Additionally, the neural network approach supports fine-tuning to incorporate new daily observations;
a significant advantage over the Monte Carlo method, which would require estimation from scratch using the new dataset.

\subsection{Implications and Future Work}
\label{subsec:implications}
The main focus of this work was to reduce arbitrage conditions by reducing the differences between the surfaces
generated with the SABR models and actual surfaces observed from daily traded option contracts.
We leveraged the ability of the implicit formulation in solving for a fixed-point problem,
combined with the very low execution times that were only possible with the Anderson acceleration method,
ensuring near optimal no-arbitrage conditions between the generated surfaces and observed implied volatility surfaces.

The approach of using neural networks, even though already significantly faster than Monte Carlo methods,
could be further optimized in the context of high-frequency intraday trading.
Adapting the feature set to include historical exogenous market factors, besides just traded contracts,
such as macroeconomic indicators, market indexes, and other financial instruments,
could allow the model to better adapt to changing market regimes and could be a promising extension in future research.
